```latex
\chapter{KẾT QUẢ THỰC NGHIỆM VÀ THẢO LUẬN}

\section{Môi trường thực nghiệm}

\subsection{Cấu hình phần cứng}

Các thí nghiệm được thực hiện trên máy tính có cấu hình như sau:

\begin{table}[H]
\centering
\caption{Cấu hình môi trường thực nghiệm}
\begin{tabular}{|l|l|}
\hline
CPU & Apple M-series / Intel Core i7 \\
\hline
RAM & 16 GB \\
\hline
Ổ cứng & SSD 512 GB \\
\hline
Hệ điều hành & macOS / Ubuntu \\
\hline
\end{tabular}
\end{table}

\subsection{Môi trường phần mềm}

\begin{table}[H]
\centering
\caption{Các công cụ sử dụng trong nghiên cứu}
\begin{tabular}{|l|l|}
\hline
Ngôn ngữ & Python 3.x \\
\hline
Pandas & Xử lý dữ liệu \\
\hline
NumPy & Tính toán số học \\
\hline
Scikit-learn & Machine Learning \\
\hline
XGBoost & Gradient Boosting \\
\hline
LightGBM & Gradient Boosting \\
\hline
CatBoost & Gradient Boosting \\
\hline
SHAP & Explainable AI \\
\hline
Matplotlib & Visualization \\
\hline
Seaborn & Visualization \\
\hline
\end{tabular}
\end{table}

\section{Phân tích dữ liệu sau tiền xử lý}

\subsection{Phân bố dữ liệu}

Sau quá trình làm sạch dữ liệu, tập dữ liệu cuối cùng bao gồm:

\[
N = 30012
\]

quan sát.

Phân bố biến mục tiêu được thể hiện trong Bảng \ref{tab:class_distribution}.

\begin{table}[H]
\centering
\caption{Phân bố các nhóm tín dụng}
\label{tab:class_distribution}
\begin{tabular}{|c|c|c|}
\hline
Nhóm & Số lượng & Tỷ lệ (\%) \\
\hline
1 & xxx & xx.xx \\
2 & xxx & xx.xx \\
3 & xxx & xx.xx \\
4 & xxx & xx.xx \\
5 & xxx & xx.xx \\
\hline
\end{tabular}
\end{table}

\begin{figure}[H]
\centering
% \includegraphics[width=0.7\textwidth]{figures/class_distribution.png}
\caption{Phân bố biến mục tiêu}
\label{fig:class_distribution}
\end{figure}

Kết quả cho thấy dữ liệu có hiện tượng mất cân bằng giữa các lớp, trong đó nhóm nợ đủ tiêu chuẩn chiếm tỷ trọng lớn nhất.

\subsection{Phân tích tương quan}

Ma trận tương quan Pearson được sử dụng để đánh giá mối quan hệ giữa các biến số.

\begin{figure}[H]
\centering
% \includegraphics[width=0.9\textwidth]{figures/correlation_heatmap.png}
\caption{Ma trận tương quan giữa các biến}
\label{fig:corr}
\end{figure}

Kết quả cho thấy các biến liên quan đến lịch sử tín dụng, số ngày quá hạn và dư nợ có mức tương quan đáng kể với biến mục tiêu.

\section{Kết quả huấn luyện mô hình}

\subsection{Logistic Regression}

Logistic Regression được sử dụng làm mô hình cơ sở (Baseline Model) do tính đơn giản và khả năng diễn giải cao \cite{hosmer2013applied}.

Bảng \ref{tab:lr_result} trình bày kết quả của mô hình.

\begin{table}[H]
\centering
\caption{Kết quả Logistic Regression}
\label{tab:lr_result}
\begin{tabular}{|c|c|}
\hline
Accuracy & xx.xx \\
\hline
Precision & xx.xx \\
\hline
Recall & xx.xx \\
\hline
F1-score & xx.xx \\
\hline
\end{tabular}
\end{table}

\subsection{Decision Tree}

Decision Tree được xây dựng nhằm khai thác các quan hệ phi tuyến trong dữ liệu \cite{quinlan1986induction}.

\begin{table}[H]
\centering
\caption{Kết quả Decision Tree}
\begin{tabular}{|c|c|}
\hline
Accuracy & xx.xx \\
\hline
Precision & xx.xx \\
\hline
Recall & xx.xx \\
\hline
F1-score & xx.xx \\
\hline
\end{tabular}
\end{table}

\subsection{Random Forest}

Random Forest giúp giảm hiện tượng overfitting bằng cách kết hợp nhiều cây quyết định \cite{breiman2001random}.

\begin{table}[H]
\centering
\caption{Kết quả Random Forest}
\label{tab:rf_result}
\begin{tabular}{|c|c|}
\hline
Accuracy & xx.xx \\
\hline
Precision & xx.xx \\
\hline
Recall & xx.xx \\
\hline
F1-score & xx.xx \\
\hline
\end{tabular}
\end{table}

\subsection{XGBoost}

XGBoost là thuật toán boosting được đánh giá rất cao trong các bài toán tín dụng \cite{chen2016xgboost}.

\begin{table}[H]
\centering
\caption{Kết quả XGBoost}
\label{tab:xgb_result}
\begin{tabular}{|c|c|}
\hline
Accuracy & xx.xx \\
\hline
Precision & xx.xx \\
\hline
Recall & xx.xx \\
\hline
F1-score & xx.xx \\
\hline
\end{tabular}
\end{table}

\subsection{LightGBM}

LightGBM được thiết kế nhằm tối ưu tốc độ huấn luyện trên dữ liệu lớn \cite{ke2017lightgbm}.

\begin{table}[H]
\centering
\caption{Kết quả LightGBM}
\label{tab:lgb_result}
\begin{tabular}{|c|c|}
\hline
Accuracy & xx.xx \\
\hline
Precision & xx.xx \\
\hline
Recall & xx.xx \\
\hline
F1-score & xx.xx \\
\hline
\end{tabular}
\end{table}

\subsection{CatBoost}

CatBoost xử lý tốt các biến phân loại và hạn chế overfitting \cite{prokhorenkova2018catboost}.

\begin{table}[H]
\centering
\caption{Kết quả CatBoost}
\label{tab:cat_result}
\begin{tabular}{|c|c|}
\hline
Accuracy & xx.xx \\
\hline
Precision & xx.xx \\
\hline
Recall & xx.xx \\
\hline
F1-score & xx.xx \\
\hline
\end{tabular}
\end{table}

\section{So sánh các mô hình}

Bảng \ref{tab:model_compare} trình bày kết quả tổng hợp.

\begin{table}[H]
\centering
\caption{So sánh các mô hình}
\label{tab:model_compare}
\begin{tabular}{|l|c|c|c|c|}
\hline
Mô hình & Accuracy & Precision & Recall & F1-score \\
\hline
Logistic Regression & xx.xx & xx.xx & xx.xx & xx.xx \\
Decision Tree & xx.xx & xx.xx & xx.xx & xx.xx \\
Random Forest & xx.xx & xx.xx & xx.xx & xx.xx \\
XGBoost & xx.xx & xx.xx & xx.xx & xx.xx \\
LightGBM & xx.xx & xx.xx & xx.xx & xx.xx \\
CatBoost & xx.xx & xx.xx & xx.xx & xx.xx \\
\hline
\end{tabular}
\end{table}

\begin{figure}[H]
\centering
% \includegraphics[width=0.9\textwidth]{figures/model_comparison.png}
\caption{So sánh Accuracy giữa các mô hình}
\end{figure}

Kết quả cho thấy các mô hình Ensemble Learning có hiệu quả cao hơn đáng kể so với Logistic Regression và Decision Tree.

\section{Đánh giá chi tiết mô hình tốt nhất}

Giả sử CatBoost đạt kết quả tốt nhất.

\subsection{Confusion Matrix}

\begin{figure}[H]
\centering
% \includegraphics[width=0.8\textwidth]{figures/catboost_confusion_matrix.png}
\caption{Confusion Matrix của CatBoost}
\label{fig:cm}
\end{figure}

Ma trận nhầm lẫn cho thấy phần lớn các khách hàng thuộc nhóm nợ đủ tiêu chuẩn được nhận diện chính xác.

Tuy nhiên các nhóm 3, 4 và 5 vẫn tồn tại hiện tượng nhầm lẫn do số lượng quan sát thấp hơn.

\subsection{ROC-AUC}

Đường cong ROC được sử dụng để đánh giá khả năng phân biệt giữa các lớp \cite{fawcett2006roc}.

\begin{figure}[H]
\centering
% \includegraphics[width=0.8\textwidth]{figures/roc_curve.png}
\caption{ROC Curve của mô hình CatBoost}
\end{figure}

Giá trị AUC trung bình đạt:

\[
AUC = xx.xx
\]

cho thấy khả năng phân loại tốt của mô hình.

\section{Giải thích mô hình bằng SHAP}

Một trong những hạn chế của các mô hình Ensemble là khả năng giải thích thấp.

Để khắc phục vấn đề này, nghiên cứu sử dụng SHAP (SHapley Additive exPlanations) \cite{lundberg2017unified}.

SHAP được xây dựng dựa trên lý thuyết giá trị Shapley:

\[
\phi_i
=
\sum_{S \subseteq N \setminus \{i\}}
\frac{|S|!(|N|-|S|-1)!}
{|N|!}
\left[
f(S \cup \{i\}) - f(S)
\right]
\]

\subsection{Feature Importance}

\begin{figure}[H]
\centering
% \includegraphics[width=0.9\textwidth]{figures/shap_importance.png}
\caption{Tầm quan trọng của các đặc trưng theo SHAP}
\end{figure}

Kết quả cho thấy các biến ảnh hưởng lớn nhất đến xếp hạng tín dụng bao gồm:

\begin{itemize}
    \item Số ngày quá hạn.
    \item Dư nợ hiện tại.
    \item Lịch sử tín dụng.
    \item Số khoản vay đang hoạt động.
    \item Tuổi khoản vay.
\end{itemize}

\subsection{SHAP Summary Plot}

\begin{figure}[H]
\centering
% \includegraphics[width=0.9\textwidth]{figures/shap_summary.png}
\caption{SHAP Summary Plot}
\end{figure}

SHAP Summary Plot cho thấy chiều hướng tác động của từng biến lên xác suất rơi vào nhóm nợ xấu.

\section{Thảo luận kết quả}

Kết quả nghiên cứu cho thấy các thuật toán Machine Learning hiện đại có khả năng hỗ trợ hiệu quả cho bài toán xếp hạng tín dụng.

Các mô hình boosting như XGBoost, LightGBM và CatBoost đạt hiệu năng cao hơn các mô hình truyền thống. Điều này phù hợp với các nghiên cứu trước đây của Baesens và cộng sự \cite{baesens2003benchmarking} cũng như Lessmann và cộng sự \cite{lessmann2015benchmarking}.

Việc áp dụng SMOTE giúp cải thiện khả năng nhận diện các nhóm nợ có tỷ lệ xuất hiện thấp. Đồng thời, SHAP giúp nâng cao khả năng giải thích của mô hình, đáp ứng yêu cầu minh bạch trong lĩnh vực tài chính ngân hàng.

Từ góc độ thực tiễn, mô hình đề xuất có thể được tích hợp vào hệ thống xếp hạng tín dụng nội bộ nhằm hỗ trợ cán bộ tín dụng trong quá trình đánh giá khách hàng và cảnh báo sớm rủi ro tín dụng.

\section{Tóm tắt chương}

Chương này đã trình bày kết quả thực nghiệm của các mô hình học máy trên bộ dữ liệu tín dụng thực tế của một ngân hàng thương mại tại Việt Nam.

Kết quả cho thấy các mô hình Ensemble Learning, đặc biệt là CatBoost và XGBoost, đạt hiệu quả dự báo cao hơn so với các mô hình truyền thống. Bên cạnh đó, phương pháp SHAP giúp giải thích các quyết định của mô hình và xác định các yếu tố ảnh hưởng quan trọng đến kết quả xếp hạng tín dụng.

Các kết quả này là cơ sở để đưa ra kết luận và khuyến nghị trong Chương 5.
```
